\documentclass{article}
\usepackage{graphicx,datetime,hyperref,verbatim,caption,subcaption} % Required for inserting images

\usepackage{geometry}
\geometry{
a4paper,
total={170mm,257mm},
left=20mm,
top=20mm,
}

\newcommand{\setlasteditor}[1]{\gdef\lasteditor{#1}}
\newcommand{\lastedited}{%
    \vspace{1mm} {\footnotesize Last edited by: \lasteditor} \vspace{3mm}
    \newline

}

\title{Formalities}
\date{September 2024}

\begin{document}

\maketitle


\tableofcontents
\newpage
\section*{Introduction}
\setlasteditor{Aayush Desai}
\lastedited
\noindent
Hello world. This is the ISTA Formalities Page, which will hopefully have all the neat tools and tricks the department will have developed over the years. The repository/document is maintained by Aayush Desai for the time being, and is open to contributions from anyone who would like to add to it.
\newline
\noindent
This document will contain the immediate formalities you need to take care of after you arrive at IST. While there will be a lot of overlap between the HR meetings that take care of these topics, this is a one-place repository of \emph{to-do} formalities at IST. Note, you must be connected to ISTA network (either ISTA wifi or VPN) to access some of the links. Please access the IT/software documentation in case you need help with that. 
% Here you can find access to example scripts that are group agnostic (mostly). Should you have any ideas for further developments, just add another section and start writing.
% \newline
% \newline
% For those working on Overleaf, once you are done with the edits, please make sure to push the changes to the repository. You can do this by:

% \begin{enumerate}
%     \item \textbf{Committing the changes:} This can be done by clicking on the "menu" button, on the top left corner of the screen, navigating to "github" and selecting commit changes. This will save the changes you have made to the repository.
%     \item \textbf{Pushing the changes:} This can be done by clicking on the three dots on the top right corner of the screen, and selecting push changes. This will push the changes you have made to the repository.
%     \item \textbf{Syncing the changes:} This can be done by clicking on the three dots on the top right corner of the screen, and selecting sync changes. This will sync the changes you have made to the repository.
% \end{enumerate}
% \noindent
% For those working on the local machine, please make sure to:

% \begin{enumerate}
%     \item \textbf{Pull from the main repo} git pull from the repository to get the latest changes.
%     \item \textbf{Make changes} make the changes you need to make.
%     \item \textbf{Push to the main repo} git add, git commit, and git push the changes to the repository.
% \end{enumerate}
\section{A2Ps}
A2Ps are sort of the most worldly-wise members of the group. They handle the cost-center, conference bookings (tickets, hotels etc), amongst many other responsibilities. They are extremely approachable and very happy to help. Table. \ref{tab:A2Ps} contains the details of the A2Ps responsible for each group. They are your first points of contact for most practical matters at IST. While one should try to reach out to the specific A2P in question but, if unavailable, the others would be happy to help!
\begin{table}[h]
    \centering
    \begin{tabular}{c|c|c}
       A2P  & Groups & Contact \\
       \hline
        Viktoriya Momchilova  & Caiazzo, Haimann and Götberg Groups & viktoriya.momchilova@ist.ac.at\\
        Benjamin Jones & Mathee Group & benjamin.jones@ist.ac.at\\
        Adrienn Keri & Bugnet Group & adrienn.keri@ist.ac.at \\
        \hline
        
    \end{tabular}
    \caption{List containing the contact details for A2Ps responsible for each Astro Group at IST. }
    \label{tab:A2Ps}
\end{table}
\section{VISA/HR Stuff}
\setlasteditor{Aayush Desai}
\lastedited
\noindent
While you would ideally have figured out the Visa-D by the time you peruse this page, please feel free to reach out to the A2Ps and Vlad Cozac (email: vlad.cozac@ist.ac.at) in case you have questions/problems.
\\
Once you arrive at ISTA, you must apply for your Meldezettel. This is essentially telling the local municipality that you have arrived and need their attention. This applies to all incoming personnel. You must apply within three days (although it's okay if you do it within a week at most) to Klosterneuburg (info \href{https://intranet.ista.ac.at/file/file/download?guid=f65ffda4-ab40-4a3c-965b-10dba02f1951\&hash\_sha1=fbf9b0f3}{here}; form \href{https://intranet.ista.ac.at/file/file/download?guid=00ca3b52-29f6-4813-962b-cc852358e9ca\&hash\_sha1=953ce7de}{here}) or Vienna (info \href{}{here}; form \href{https://intranet.ista.ac.at/file/file/download?guid=00ca3b52-29f6-4813-962b-cc852358e9ca\&hash\_sha1=953ce7de}{here}), depending on where you stay. 
\\
\noindent
Depending on your duration of stay and country of origin, you may have to apply for a long-term residence permit. Most of the information is available \href{https://hr.pages.ist.ac.at/residence-permit-registration-in-austria/}{here}.
\subsection{Emergency Visa to leave the country}
\noindent
In case you want to leave the country while you are still awaiting your (decision on) residence permit, you can directly apply for an emergency \emph{nonvignette} visa (info \href{https://www.wien.gv.at/amtswege/notvignette}{here}). The usual processing time is 2 days from the date of application (which is just an email) and you can collect the stamped visa. The cost for this is roughly \textbf{45 Eur}.

\subsection{Ecard}
\noindent
While you are insured as soon as you start your employment at IST, it is important that you apply for an \emph{ECard}. It is sort of an all-powerful card that contains all your health records (sus?) from your first visit to a doctor. Information for the Ecard can be found \href{https://intranet.ista.ac.at/file/file/download?guid=a4e5b7c0-6074-425d-ad33-fa05696081cb&hash_sha1=0ec03d01}{here}: 
\subsection{Passport sized photo}
\noindent
For most of your Austrian document applications, you will have to submit a passport-sized photo, with the date it was captured on printed. One can easily get such a photo at one of the photo booths at Spittleau station (\href{https://maps.app.goo.gl/CRdGfkwThCaowHny5}{here}) or Heiligenstadt station (\href{https://maps.app.goo.gl/hUrvxbjUpanRm38Y7}{here}) or at Kierling station (\href{https://maps.app.goo.gl/9DMVMN9KAu2Senhv6}{here}). The photo booths are usually small shacks, located within the station. It usually costs around \textbf{10 Eur} and they accept card (except Heiligenstadt which only accepts cash). The date is automatically printed on the printed photos. Please do not cut the photos yourself and instead take the full mosaic to your appointment. 
\section{Slack Channels to be on}
\setlasteditor{Aayush Desai}
\lastedited
\noindent
This is a list of Slack channels that one is recommended to join on the \textit{ISTAstro} Slack workspace, to keep abreast with all the latest developments. Just click the link to join these public channels. (This list will be updated as and when new public channels are created)
\begin{enumerate}
    \item General: \href{}{https://istastro.slack.com/archives/C059YAJN8KS}
    \item Astro Student Sessions: \href{https://istastro.slack.com/archives/C062NF7FAB1}{link}
    \item Astro Seminars: \href{https://istastro.slack.com/archives/C06JJNEFFPY}{link}
    \item Astro Student Sessions: \href{https://istastro.slack.com/archives/C062NF7FAB1}{link}
    \item Astro Classes: \href{https://istastro.slack.com/archives/C08T3N94363}{link}
    \item Astro Conferences: \href{https://istastro.slack.com/archives/C06628E3ELR}{link}
    \item Food: \href{https://istastro.slack.com/archives/C05V5H49W3A}{link}
    \item HPC-Computing: \href{https://istastro.slack.com/archives/C07RQ876FC6}{link}
    \item Long Night of Research: \href{https://istastro.slack.com/archives/C06UYT9NZ4P}{link}
    \item Lunch/Dinner with Visitors: \href{https://istastro.slack.com/archives/C06ECTPKPK7}{link}
    \item Palomar: \href{https://istastro.slack.com/archives/C09G29VSU7L}{link}
    \item Random: \href{https://istastro.slack.com/archives/C05QQ8KRG7R}{link}
    \item Social: \href{https://istastro.slack.com/archives/C05R4QFL66N}{link}
    \item Vista Science Center: \href{https://istastro.slack.com/archives/C089Z3R6UF4}{link}
\end{enumerate}



\end{document}
